\documentclass[a4paper,12pt]{article}

\usepackage[margin=2.5cm]{geometry}
\usepackage{graphicx}
\usepackage{natbib}
\usepackage{hyperref}
\usepackage{amsmath}
\usepackage{amssymb}
\usepackage{booktabs}
\usepackage{longtable}
\usepackage{placeins}
\usepackage{listings}
\usepackage{xcolor}
\usepackage{url}

\lstset{
  basicstyle=\ttfamily\small,
  breaklines=true,
  breakatwhitespace=false,
  columns=fullflexible,
  showspaces=false,
  frame=single,
  numbers=left,
  numberstyle=\tiny,
}

\makeatletter
\def\lst@visiblespace{\ }
\makeatother

\title{Classification of Emission-Line Galaxies with SDSS and WISE}
\date{\today}
\author{Zain Ul Abideen\\[2pt]
  \small\textit{Consortium of the Erasmus Mundus Joint Master Programme}\\
  \small\textit{in Astrophysics and Space Science (MASS)}\\[2pt]
  \small\texttt{zain.ulabideen@students.uniroma2.eu}\\[2pt]
  \small\url{www.master-mass.eu}}

\begin{document}
% Auto-generated by scripts/07_generate_numbers.py -- do not edit
\newcommand{\samplesize}{9998}
\newcommand{\bptSfCount}{6285}
\newcommand{\bptCompCount}{2277}
\newcommand{\bptAgnCount}{1436}
\newcommand{\bptSfPct}{62.9}
\newcommand{\bptCompPct}{22.8}
\newcommand{\bptAgnPct}{14.4}
\newcommand{\kauffSfCount}{6285}
\newcommand{\kauffAgnCount}{3713}
\newcommand{\kauffAgnPct}{37.1}
\newcommand{\whanSfCount}{5117}
\newcommand{\whanSfPct}{51.2}
\newcommand{\whanStrongCount}{2810}
\newcommand{\whanStrongPct}{28.1}
\newcommand{\whanWagnCount}{853}
\newcommand{\whanWagnPct}{8.5}
\newcommand{\whanRetCount}{1207}
\newcommand{\whanRetPct}{12.1}
\newcommand{\whanPassCount}{11}
\newcommand{\whanPassPct}{0.1}
\newcommand{\whanAgnTotCount}{3663}
\newcommand{\whanAgnTotPct}{36.6}
\newcommand{\bothAgnCount}{673}
\newcommand{\bptOnlyCount}{763}
\newcommand{\whanOnlyCount}{2990}
\newcommand{\bothAgnPct}{6.7}
\newcommand{\bptOnlyPct}{7.6}
\newcommand{\whanOnlyPct}{29.9}
\newcommand{\woStrongCount}{2391}
\newcommand{\woWagnCount}{599}
\newcommand{\woStrongPct}{80.0}
\newcommand{\woWagnPct}{20.0}
\newcommand{\woBptSfCount}{1207}
\newcommand{\woBptCompCount}{1783}
\newcommand{\woBptSfPct}{40.4}
\newcommand{\woBptCompPct}{59.6}
\newcommand{\boRetCount}{746}
\newcommand{\boSfCount}{9}
\newcommand{\boPassCount}{8}
\newcommand{\boRetPct}{97.8}
\newcommand{\boSfPct}{1.2}
\newcommand{\boPassPct}{1.0}
\newcommand{\wiseAgnSample}{1436}
\newcommand{\wiseMatched}{1408}
\newcommand{\wiseMatchRate}{98.1}
\newcommand{\wiseWedge}{271}
\newcommand{\wiseWedgePct}{19.2}
\newcommand{\wiseOutside}{1137}
\newcommand{\wiseOutsidePct}{80.8}
\newcommand{\siiSample}{8718}
\newcommand{\siiAboveLimit}{3184}
\newcommand{\siiAboveLimitPct}{36.5}
\newcommand{\neValid}{5647}
\newcommand{\sigmaMedian}{78}
\newcommand{\sigmaMin}{1}
\newcommand{\sigmaMax}{415}
\newcommand{\sigmaSubTen}{11}
\newcommand{\mbhNtrem}{8697}
\newcommand{\mbhNgeb}{8704}
\newcommand{\mbhNwu}{8695}
\newcommand{\mbhMedTrem}{6.50}
\newcommand{\mbhMedGeb}{6.56}
\newcommand{\mbhMedWu}{6.43}
\newcommand{\mbhBelowFour}{21}
\newcommand{\mbhBinFiveSixCount}{2361}
\newcommand{\mbhBinFiveSixPct}{27.1}
\newcommand{\mbhBinSixSevenCount}{3583}
\newcommand{\mbhBinSixSevenPct}{41.2}
\newcommand{\mbhBinSevenEightCount}{2299}
\newcommand{\mbhBinSevenEightPct}{26.4}
\newcommand{\mbhBinEightNineCount}{308}
\newcommand{\mbhBinEightNinePct}{3.5}
\newcommand{\mbhBinOtherCount}{146}
\newcommand{\mbhBinOtherPct}{1.7}

\maketitle

\begin{abstract}
We classify $\sim$10,000 narrow emission-line galaxies at $z < 0.35$ from SDSS DR18
using optical BPT and WHAN diagnostic diagrams.  The BPT (Kewley+01) three-class
scheme gives \bptSfPct{}\% star-forming, \bptCompPct{}\% composite, and \bptAgnPct{}\% AGN/Seyfert
galaxies.  The WHAN diagram identifies a larger AGN fraction (36.6\,\%), mostly
Strong AGN missed by BPT.  We cross-match the BPT AGN with WISE and find that
1,408 of \wiseAgnSample{} (98.0\,\%) have mid-infrared counterparts; \wiseWedge{} (\wiseWedgePct{}\,\%) fall
inside the Mateos+12 AGN selection wedge.  Using the [S\,II]~$\lambda\lambda$6717,6731
flux ratio, we estimate electron densities for \neValid{} galaxies. We estimate
supermassive black hole masses via the $M_{\rm BH}$--$\sigma$ relation of
\citet{tremaine2002} and compare results with the \citet{gebhardt2000} and \citet{wu2009} calibrations.
\end{abstract}

%---------------------------------------------------------------------
\section{Introduction}
%---------------------------------------------------------------------

When a telescope takes a spectrum of a galaxy, the emission lines in that
spectrum tell us what is ionising the gas.  Hot young stars produce one
pattern of line ratios; an accreting supermassive black hole (an Active
Galactic Nucleus, or AGN) produces a different, ``harder'' pattern.
(The most common AGN in the local Universe are Seyfert galaxies, which
have bright, point-like nuclei with strong emission lines.)

The standard way to separate star-forming galaxies from AGN is the BPT
diagram \citep{baldwin1981,veilleux1987}, which plots
$\log\,[{\rm O\,III}]/{\rm H}\beta$ against
$\log\,[{\rm N\,II}]/{\rm H}\alpha$.  \citet{kewley2001} calculated the
theoretical boundary that pure starburst galaxies cannot cross (the
``maximum starburst line''), while \citet{kauffmann2003} drew an empirical
line that best separates the two populations in real data.

The BPT diagram has a limitation: it needs all four lines to be detected
with good signal-to-noise.  \citet{cidfernandes2011} introduced the WHAN
diagram, which uses only [N\,II]/H$\alpha$ and the H$\alpha$ equivalent
width.  The WHAN diagram can classify galaxies that the BPT misses, and
it can identify ``retired galaxies'' whose ionisation comes from old
stars rather than an AGN.

Mid-infrared colours from the WISE satellite \citep{wright2010} offer a
different way to find AGN.  Hot dust near the black hole glows brightly
at 12\,$\mu$m, producing a distinctive position in the
$\log(f_{12}/f_{4.6})$ vs.\ $\log(f_{4.6}/f_{3.4})$ plane.
\citet{mateos2012} defined a wedge in this plane that selects luminous
AGN with high completeness.

In this project we: (1) classify \samplesize{} SDSS galaxies with the BPT and
WHAN diagrams; (2) cross-match the AGN with WISE and apply the Mateos+12
mid-infrared selection; (3) measure the electron density of the
narrow-line region gas and estimate black hole masses; and (4) look for
trends in the gas properties across the BPT plane.

%---------------------------------------------------------------------
\section{Data and Sample Selection}
\label{sec:data}
%---------------------------------------------------------------------

\subsection{SDSS Spectroscopic Data}

We queried the Sloan Digital Sky Survey Data Release 18
SkyServer\footnote{\url{https://skyserver.sdss.org/dr18/}} using
paginated SQL requests.  The query joined the \texttt{SpecObj} table
with \texttt{galSpecLine}; in DR18 this single table provides
emission-line fluxes, equivalent widths, and velocity dispersions for
all the lines we need.

We selected galaxies with:
\begin{itemize}
  \item Redshift $0 < z < 0.35$;
  \item Spectroscopic class \texttt{'GALAXY'};
  \item Signal-to-noise ratio $\mathrm{S/N} \ge 3$ on all four
        diagnostic lines: H$\alpha$, H$\beta$,
        [O\,III]~$\lambda$5007, and [N\,II]~$\lambda$6584.
\end{itemize}

The query returned 10,000 galaxies; after removing 2 with physically
impossible line fluxes (likely fiber contamination), the final sample
contains \samplesize{} galaxies matching the criteria
(sorted by \texttt{specObjID}), which we use for all
subsequent analysis.  (The \texttt{specObjID} incorporates the
SDSS plate and MJD, so the sample is slightly biased toward
earlier-epoch SDSS observations.)  The redshift range, S/N cut, and
``GALAXY'' class requirement ensure bright, well-measured narrow
emission lines.

\subsection{WISE Cross-Matching}

We cross-matched the BPT AGN subsample (\wiseAgnSample{} galaxies) with the WISE
All-Sky catalog via the SDSS SkyServer's \texttt{wise\_allsky} table.
We used a matching radius of $3''$ (about half the WISE beam size at
3.4\,$\mu$m).  For each SDSS object, we queried WISE sources within a
bounding box around the position and kept the nearest match within the
$3''$ tolerance.

%---------------------------------------------------------------------
\section{Task 1: BPT and WHAN Diagrams}
\label{sec:task1}
%---------------------------------------------------------------------

\subsection{BPT Diagram}

We computed the two line ratios from the SDSS flux measurements:
\[
\log\,([{\rm O\,III}]/{\rm H}\beta) \quad\text{and}\quad
\log\,([{\rm N\,II}]/{\rm H}\alpha) .
\]

Figure~\ref{fig:bpt} shows the result.  Objects below the Kauffmann+03
line are star-forming; objects between the two curves are composite
(ionised by a mix of star formation and AGN activity); objects above
the Kewley+01 line are classified as AGN/Seyfert.  Objects with
$\log\,([{\rm N\,II}]/{\rm H}\alpha) \ge 0.47$---beyond the Kewley
curve's valid range---are all in the high-excitation LINER/Seyfert 
region and are classified as AGN.  (LINERs---Low-Ionization Nuclear
Emission-line Regions---are a class of AGN with weaker [O\,III] relative
to [N\,II] than classical Seyferts; our BPT analysis does not separate
these two sub-classes.)

When we collapse to the simpler two-class Kauffmann+03 scheme
(star-forming below the line, AGN above), we get \kauffSfCount{} star-forming
galaxies and \kauffAgnCount{} AGN (\kauffAgnPct{}\%).

\begin{figure}[htbp]
  \centering
  \includegraphics[width=0.55\textwidth]{../output/bpt_diagram.pdf}
  \caption{BPT diagram.  Solid line: Kewley+01 maximum starburst line.
           Dashed line: Kauffmann+03 empirical division.
           Blue: star-forming; grey: composite; red: AGN/Seyfert.}
  \label{fig:bpt}
\end{figure}

\begin{table}[htbp]
  \centering
  \caption{BPT (Kewley+01 three-class) classification.}
  \label{tab:bpt_stats}
  \begin{tabular}{lrr}
    \toprule
    Class & $N$ & Fraction \\
    \midrule
    Star-forming & \bptSfCount{} & \bptSfPct{}\,\% \\
    Composite    & \bptCompCount{} & \bptCompPct{}\,\% \\
    AGN/Seyfert  & \bptAgnCount{} & \bptAgnPct{}\,\% \\
    \bottomrule
  \end{tabular}
\end{table}

\subsection{WHAN Diagram}

The WHAN diagram uses H$\alpha$ equivalent width (EW) on the vertical
axis instead of [O\,III]/H$\beta$.  We took the absolute value of the
EW from \texttt{galSpecLine}, since DR18 stores emission-line EWs as
negative numbers.  The classification boundaries follow
\citet{cidfernandes2011}:
\begin{itemize}
  \item \textbf{Star-forming}: $\log([{\rm N\,II}]/{\rm H}\alpha) \le -0.4$
        and $W_{{\rm H}\alpha} > 3$\,\AA;
  \item \textbf{Strong AGN}: $\log([{\rm N\,II}]/{\rm H}\alpha) > -0.4$
        and $W_{{\rm H}\alpha} > 6$\,\AA;
  \item \textbf{Weak AGN (wAGN)}: $\log([{\rm N\,II}]/{\rm H}\alpha) > -0.4$
        and $3 < W_{{\rm H}\alpha} \le 6$\,\AA;
  \item \textbf{Retired galaxies}: $W_{{\rm H}\alpha} < 3$\,\AA\ (but
        $\ge 0.5$\,\AA);
  \item \textbf{Passive galaxies}: $W_{{\rm H}\alpha} < 0.5$\,\AA.
\end{itemize}

Following the task instructions, Figure~\ref{fig:whan} colours points
by their BPT classification, not by WHAN class, so we can see where
each BPT type falls in the WHAN plane.

\begin{figure}[htbp]
  \centering
  \includegraphics[width=0.55\textwidth]{../output/whan_diagram.pdf}
  \caption{WHAN diagram.  Points are coloured by BPT class (Kewley+01).
           Blue: star-forming; grey: composite; red: AGN/Seyfert.
           Dashed lines show the WHAN classification boundaries at
           $\log([{\rm N\,II}]/{\rm H}\alpha)=-0.4$ and
           $W_{{\rm H}\alpha}=3,\,6$\,\AA.}
  \label{fig:whan}
\end{figure}

\begin{table}[htbp]
  \centering
  \caption{WHAN classification (Cid Fernandes+\whanPassCount{}).}
  \label{tab:whan_stats}
  \begin{tabular}{lrr}
    \toprule
    Class & $N$ & Fraction \\
    \midrule
    Star-forming & \whanSfCount{} & \whanSfPct{}\,\% \\
    Strong AGN   & \whanStrongCount{} & \whanStrongPct{}\,\% \\
    wAGN         &   \whanWagnCount{} &  \whanWagnPct{}\,\% \\
    Retired      & \whanRetCount{} & \whanRetPct{}\,\% \\
    Passive      &    \whanPassCount{} &  \whanPassPct{}\,\% \\
    \bottomrule
  \end{tabular}
\end{table}

\subsection{Comparing BPT and WHAN}

The two diagrams give very different AGN fractions: the BPT finds
\wiseAgnSample{} AGN (14.4\,\%), while the WHAN finds \whanAgnTotCount{} when combining Strong
AGN and wAGN (36.6\,\%).  Only \bothAgnCount{} galaxies (\bothAgnPct{}\,\% of the full
sample) are AGN by both methods.

The \whanOnlyCount{} galaxies that WHAN classifies as AGN but BPT does not are
mostly Strong AGN (\woStrongCount{}, \woStrongPct{}\%), with the rest being wAGN (599,
20.1\,\%).  These objects have $\log\,([{\rm N\,II}]/{\rm H}\alpha)$ high
enough to be WHAN AGN but lie in the star-forming or composite regions
of the BPT diagram (40\,\% are BPT Star-forming, 60\,\% BPT
Composite)---they are low-excitation objects that the optical diagnostic
classifies less strictly.

The \bptOnlyCount{} BPT-only AGN (those classified as AGN/Seyfert by BPT but not
by WHAN) are almost entirely Retired galaxies by the WHAN scheme (746,
97.8\,\%), with a few Star-forming (\boSfCount{}) and Passive (\boPassCount{}).  These are
galaxies whose BPT line ratios look like AGN but whose H$\alpha$
equivalent widths are too weak ($<3$\,\AA) for WHAN to flag them as
AGN.  They likely host low-luminosity AGN or post-starburst
populations.

%---------------------------------------------------------------------
\section{Task 2: WISE Mid-Infrared Colours}
\label{sec:task2}
%---------------------------------------------------------------------

Figure~\ref{fig:wise} shows the WISE colour-colour diagram for the
1,408 BPT AGN with reliable mid-infrared photometry.  The diagram is
overlaid with the Mateos+12 AGN selection wedge.

The wedge is defined by:
\begin{align}
  y &= 0.315\,x \quad \text{(centre line)} \\
  y_{\rm top}    &= 0.315\,x + 0.297 \\
  y_{\rm bottom} &= 0.315\,x - 0.\whanPassCount{}0 \\
  y_{\rm PL}  &= -3.172\,x + 0.436 \quad (\alpha = -0.3\ \text{limit})
\end{align}
where $x = \log_{10}(f_{12\,\mu{\rm m}}/f_{4.6\,\mu{\rm m}})$ and
$y = \log_{10}(f_{4.6\,\mu{\rm m}}/f_{3.4\,\mu{\rm m}})$.
We converted WISE magnitudes to fluxes using
$f_\nu \propto 10^{-m/2.5}$; the zero-points cancel in the ratios.

\begin{figure}[htbp]
  \centering
  \includegraphics[width=0.55\textwidth]{../output/wise_diagram.pdf}
  \caption{WISE colour-colour diagram.  The grey wedge is the Mateos+12
           AGN selection region.  Red points are WISE AGN; blue points
           are outside the wedge.}
  \label{fig:wise}
\end{figure}

\begin{table}[htbp]
  \centering
  \caption{WISE cross-match and AGN classification.}
  \label{tab:wise_stats}
  \begin{tabular}{lr}
    \toprule
    Metric & Value \\
    \midrule
    BPT AGN sample & \wiseAgnSample{} \\\\
    With WISE match within $3''$ & \wiseMatched{} (\wiseMatchRate{}\,\%) \\
    With good WISE photometry & \wiseMatched{} \\\\
    In Mateos+12 wedge (WISE AGN) & \wiseWedge{} (\wiseWedgePct{}\,\% of matched) \\
    AGN by both BPT and WISE & \wiseWedge{} \\\\
    Outside wedge & \wiseOutside{} (\wiseOutsidePct{}\,\%) \\
    \bottomrule
  \end{tabular}
\end{table}

Only 19.2\,\% of the optically-selected AGN also have WISE mid-infrared
colours that place them in the Mateos+12 wedge.  This is expected: the
Mateos+12 wedge was designed to select luminous AGN with bright hot
dust, while the BPT diagram catches lower-luminosity Seyferts and
LINERs that may not have strong mid-infrared excess.  The 80.8\,\% of
BPT AGN outside the wedge are likely weaker sources---low-luminosity
Seyferts or systems where the host galaxy's starlight overwhelms
the AGN in the WISE bands.

%---------------------------------------------------------------------
\section{Task 3: Electron Density, Kinematics, and Black Hole Mass}
\label{sec:task3}
%---------------------------------------------------------------------

\subsection{\texorpdfstring{Electron Density from [S\,II]}{Electron Density from [S II]}}

The ratio of the two [S\,II] lines depends on how densely packed the
electrons are in the gas.  When the electron density $n_e$ is low
($\lesssim 10^2$\,cm$^{-3}$), the 6717\,\AA\ line is stronger than the
6731\,\AA\ line (ratio $\approx 1.4$).  At higher densities, collisions
redistribute the electrons between the two energy levels, and the ratio
drops toward $\sim 0.45$.

We computed $n_e$ using the PyNeb \texttt{temden} routine,
following the methodology of \citet{sanders2016}, with an assumed electron temperature
$T_e = 10,000$\,K, typical for the narrow-line region.  (The [S\,II]
ratio depends on both density and temperature; without an independent
$T_e$ measurement the derived densities carry a systematic uncertainty
of about 0.2--0.3\,dex.)

Of the 10,000 originally queried galaxies, 8,720 had S/N $\ge 3$ on both
[S\,II] lines; after removing the 2 contaminated spectra, \siiSample{} remain.
PyNeb returned valid densities for \neValid{} of these---the rest have
ratios outside the theoretical range of the [S\,II] diagnostic.
(Approximately \siiAboveLimitPct{}\% of all measured ratios exceed the strict
low-density limit of $\sim$1.45, though PyNeb's internal convergence
criteria accept a slightly wider range, rejecting 35.2\,\%; in either
case, noise in the faint [S\,II] lines inflates many ratios beyond the
physically meaningful range.)

Figure~\ref{fig:ne_bpt} shows the median electron density across the
BPT plane.  Most galaxies have $n_e$ between 10 and $10^3$\,cm$^{-3}$,
with no strong gradient across the classification boundaries.

\begin{figure}[htbp]
  \centering
  \includegraphics[width=0.55\textwidth]{../output/bpt_density_map.pdf}
  \caption{Electron density map on the BPT plane.  Colour scale:
           $\log_{10}(n_e/{\rm cm}^{-3})$.}
  \label{fig:ne_bpt}
\end{figure}

\subsection{\texorpdfstring{[O\,III] Velocity Dispersion}{[O III] Velocity Dispersion}}

The [O\,III] and [S\,II] lines are forbidden transitions that only form
in low-density gas ($n_e \lesssim 10^4$\,cm$^{-3}$); at the much higher
densities of the Broad-Line Region these lines are collisionally
suppressed.  This identifies the measured gas with the Narrow-Line
Region.  We use \texttt{sigma\_forbidden}, which measures only the
forbidden-line widths, further isolating the NLR kinematics from any
broad Balmer-line component.

We used the \texttt{sigma\_forbidden} column from \texttt{galSpecLine},
which gives the velocity dispersion of the forbidden emission lines
(after instrumental correction).  All \siiSample{} objects with good [S\,II]
measurements also have valid $\sigma_{[{\rm O\,III}]}$ values (range
\sigmaMin{}--\sigmaMax{}~km\,s$^{-1}$---the 8 objects capped at exactly 500~km\,s$^{-1}$
by the SDSS pipeline are in the full 9,998-galaxy sample but fall
outside the S\,II S/N cut).  The median is 78~km\,s$^{-1}$.
\sigmaSubTen{} objects have $\sigma < 10$~km\,s$^{-1}$, well below the SDSS
instrumental resolution of $\sim$70~km\,s$^{-1}$; these are dominated
by the instrumental correction and should be treated with caution.

Figure~\ref{fig:sigma_bpt} maps $\sigma_{[{\rm O\,III}]}$ across the
BPT diagram.  Velocity dispersions are generally higher in the AGN
region of the diagram, which may reflect a combination of deeper
gravitational potentials and AGN-driven gas motions.  We note that
$\sigma_{[{\rm O\,III}]}$ correlates with redshift ($r \approx 0.41$
across the full sample) due to the 3$''$ SDSS fiber covering a larger
physical area at higher $z$, sampling more of the galaxy's rotation;
comparisons of $\sigma$ across redshifts should account for this
aperture bias.  Additionally, 8 objects have $\sigma_{[{\rm O\,III}]}$
capped at exactly 500~km\,s$^{-1}$ by the SDSS pipeline and should be
treated as lower limits (these are in the parent sample but not in the
S\,II-filtered subset).

\begin{figure}[htbp]
  \centering
  \includegraphics[width=0.55\textwidth]{../output/bpt_sigma_map.pdf}
  \caption{[O\,III] velocity dispersion map on the BPT plane.}
  \label{fig:sigma_bpt}
\end{figure}

\subsection{Supermassive Black Hole Mass}

Supermassive black hole mass correlates with the velocity dispersion of
the stars (and gas) in the host galaxy's bulge---the $M_{\rm
BH}$--$\sigma$ relation \citep{tremaine2002}.  We used:
\begin{equation}
  \log\left(\frac{M_{\rm BH}}{M_\odot}\right)
  = 8.13 + 4.02\,\log\left(\frac{\sigma}{200\ {\rm km\,s}^{-1}}\right) .
\end{equation}
The normalisation point 200~km\,s$^{-1}$ corresponds to the velocity
dispersion of a typical massive galaxy bulge.

We clipped velocity dispersions to 10--500~km\,s$^{-1}$ before applying
the formula to avoid unphysical results from measurement outliers.
(Note that the lower clip of 10~km\,s$^{-1}$ is below the SDSS instrumental
resolution of $\sim$70~km\,s$^{-1}$; values this low are dominated by
the instrumental correction but are retained to avoid biasing the
sample against low-mass galaxies.)
Of the 8,718 galaxies with $\sigma_{[{\rm O\,III}]}$ measurements,
\mbhNtrem{} give $M_{\rm BH}$ estimates within $10^4$--$10^{10}\,M_\odot$
(\mbhBelowFour{} objects have $\log M_{\rm BH} \le 4$ and are excluded from the
primary bins; the remaining \mbhBinOtherCount{} fall in the catch-all bin
`Below $10^5$ or above $10^9$' shown in Table~\ref{tab:mbh_bins}).

\textbf{Important caveat:} The \citet{tremaine2002} relation was calibrated on
\emph{stellar} velocity dispersions.  We are applying it to
\emph{gas} kinematics ($\sigma_{[{\rm O\,III}]}$), which can be
affected by AGN-driven outflows.  The black hole masses derived here
should be treated as indicative proxies, not precise measurements.

\begin{table}[htbp]
  \centering
  \caption{SMBH mass distribution (\citet{tremaine2002}) and comparison with
           \citet{gebhardt2000} and \citet{wu2009}.}
  \label{tab:mbh_bins}
  \begin{tabular}{p{4.5cm}rr}
    \toprule
    \textbf{Panel A:} Mass bins (\citealp{tremaine2002}) & $N$ & Fraction \\
    \midrule
    $10^5$--$10^6\,M_\odot$   & \mbhBinFiveSixCount{} & \mbhBinFiveSixPct{}\,\% \\
    $10^6$--$10^7\,M_\odot$   & \mbhBinSixSevenCount{} & \mbhBinSixSevenPct{}\,\% \\
    $10^7$--$10^8\,M_\odot$   & \mbhBinSevenEightCount{} & \mbhBinSevenEightPct{}\,\% \\
    $10^8$--$10^9\,M_\odot$   &   \mbhBinEightNineCount{} &  \mbhBinEightNinePct{}\,\% \\
    Below $10^5$ or above $10^9$ & \mbhBinOtherCount{} &  \mbhBinOtherPct{}\,\% \\
    \midrule
    \textbf{Panel B:} Relation & Median $\log(M_{\rm BH}/M_\odot)$ & $N$ \\
    \midrule
    \citealp{tremaine2002} & \mbhMedTrem{} & \mbhNtrem{} \\
    \citealp{gebhardt2000} & \mbhMedGeb{} & \mbhNgeb{} \\
    \citealp{wu2009}     & \mbhMedWu{} & \mbhNwu{} \\
    \bottomrule
  \end{tabular}
\end{table}

Table~\ref{tab:mbh_bins} (Panel A) shows the mass distribution from our
primary \citet{tremaine2002} estimate.  The sample peaks in the
$10^6$--$10^7\,M_\odot$ bin (4\boSfPct{}\%), which is typical for galaxies
at $z < 0.35$ in the SDSS spectroscopic sample.  Panel B compares the
three calibrations; the slightly different $N$ values reflect the
different slopes of each relation, which shift a few objects across
the $\log M_{\rm BH} = 4$ boundary.  All three give consistent median
masses of $3$--$4 \times 10^6\,M_\odot$.

Figure~\ref{fig:mbh_hist} shows the full distribution and the split by
BPT class.  AGN/Seyfert galaxies tend toward slightly higher black hole
masses than star-forming galaxies, consistent with the idea that more
massive black holes can drive higher ionisation states.

\begin{figure}[htbp]
  \centering
  \includegraphics[width=0.48\textwidth]{../output/mbh_histogram.pdf}
  \includegraphics[width=0.48\textwidth]{../output/mbh_by_class_histogram.pdf}
  \caption{Left: SMBH mass histogram (\citealp{tremaine2002}).  Right: mass
           distributions split by BPT class.}
  \label{fig:mbh_hist}
\end{figure}

\subsection{Kinematics Binned by Black Hole Mass}

We divided the sample into 1-dex bins in $M_{\rm BH}$ and plotted the
[O\,III] velocity dispersion maps on the BPT plane for each bin.  If
the NLR kinematics were dominated by outflows from the central engine,
we might expect higher dispersions in more massive black holes.
The binned maps are shown in Figure~\ref{fig:sigma_mbh}.

\begin{figure}[htbp]
  \centering
  \includegraphics[width=0.48\textwidth]{../output/bpt_sigma_mbh_5_6.pdf}\hfill
  \includegraphics[width=0.48\textwidth]{../output/bpt_sigma_mbh_6_7.pdf}\\[4pt]
  \includegraphics[width=0.48\textwidth]{../output/bpt_sigma_mbh_7_8.pdf}\hfill
  \includegraphics[width=0.48\textwidth]{../output/bpt_sigma_mbh_8_9.pdf}
  \caption{[O\,III] velocity dispersion on the BPT plane, binned by
           $M_{\rm BH}$ (Tremaine+02).  Top-left: $10^5$--$10^6\,M_\odot$
           (\mbhBinFiveSixCount{} objects).  Top-right: $10^6$--$10^7\,M_\odot$
           (\mbhBinSixSevenCount{} objects).  Bottom-left: $10^7$--$10^8\,M_\odot$
           (\mbhBinSevenEightCount{} objects).  Bottom-right: $10^8$--$10^9\,M_\odot$
           (\mbhBinEightNineCount{} objects).}
  \label{fig:sigma_mbh}
\end{figure}

The most massive bin ($10^8$--$10^9\,M_\odot$, 308 objects) shows a
median $\sigma_{[{\rm O\,III}]}$ of 168~km\,s$^{-1}$, compared to
\sigmaMedian{}~km\,s$^{-1}$ in the $10^5$--$10^6\,M_\odot$ bin---an increase
consistent with the M-sigma relation itself rather than evidence for
AGN-driven outflows.  Most galaxies across all mass bins have
$\sigma_{[{\rm O\,III}]} \sim 50$--150~km\,s$^{-1}$, consistent with
gravitational motions in the host bulge rather than strong outflows.

%---------------------------------------------------------------------
\section{Discussion}
\label{sec:discussion}
%---------------------------------------------------------------------

\subsection{Why BPT and WHAN Disagree}

The large difference between the BPT AGN fraction (14.4\,\%) and the
WHAN AGN fraction (36.6\,\%) comes from the different information the
two diagrams use.  BPT compares two line ratios; WHAN replaces the
[O\,III]/H$\beta$ ratio with the H$\alpha$ equivalent width.  An object
can have AGN-like [N\,II]/H$\alpha$ while having [O\,III]/H$\beta$
too low to cross the Kauffmann line---the WHAN diagram catches these
objects because it does not require strong [O\,III].

The \woStrongCount{} WHAN-only Strong AGN are exactly this population:
high-[N\,II] objects whose [O\,III] emission is weak enough to place
them mostly in the composite region of the BPT diagram (60\,\%), with
the rest in the star-forming region (40\,\%).  They may be
low-luminosity AGN or systems where the narrow-line region gas is
distant from the ionising source.

The 746 BPT-only AGN that WHAN calls Retired have strong enough
[O\,III] to reach the Seyfert region, but their H$\alpha$ equivalent
widths are small ($<3$\,\AA).  This is the signature of
post-starburst or weakly accreting systems---the line ratios still look
like an AGN, but the total line flux relative to the continuum is low.

\subsection{Optical vs.\ Mid-Infrared AGN}

The Mateos+12 WISE wedge identifies only 19.2\,\% of our BPT AGN as
mid-infrared AGN.  This is a feature, not a bug: the wedge selects
luminous, unobscured AGN whose hot dust dominates the mid-infrared
colours.  The 80.8\,\% of BPT AGN outside the wedge are likely weaker
sources---low-luminosity Seyferts or systems where the host galaxy's
stellar light overwhelms the AGN in the WISE bands.  A cross-match of
the full sample (not just the BPT AGN) might reveal
WISE-selected AGN that are optically classified as star-forming or
composite.

\subsection{NLR Gas Properties}

The electron densities we measure ($n_e \sim 10$--$10^3$\,cm$^{-3}$)
are typical of narrow-line region gas.  The lack of a strong gradient
across the BPT plane suggests that the local gas density is not the
main parameter controlling whether a galaxy appears as star-forming,
composite, or AGN in the BPT diagram.  The [O\,III] velocity
dispersions are also fairly uniform, although the highest values appear
in the AGN region, which may indicate some contribution from
non-gravitational motions.

\subsection{Black Hole Masses}

The SMBH mass distribution peaks at $M_{\rm BH} \sim 3 \times
10^6\,M_\odot$, matching expectations for the typical SDSS galaxy at
$z < 0.35$.  The three M-sigma calibrations agree within $\sim 0.13$\,dex
in the median, which is well within the intrinsic scatter of the
relation ($\sim 0.3$\,dex).  The tendency for AGN/Seyfert galaxies to
have higher $M_{\rm BH}$ than star-forming galaxies is consistent with
more massive black holes being able to ionise a larger fraction of the
host galaxy's gas.

%---------------------------------------------------------------------
\section{Conclusions}
\label{sec:conclusions}

We analysed \samplesize{} narrow emission-line galaxies from SDSS DR18.  Our
main results are:

\begin{enumerate}
  \item The BPT diagram classifies \bptSfPct{}\% as star-forming,
        \bptCompPct{}\% as composite, and \bptAgnPct{}\% as AGN/Seyfert.
  \item The WHAN diagram finds a much larger AGN fraction
        (\whanAgnTotPct{}\%), driven by Strong AGN that fall on the
        star-forming side of the BPT diagram.
  \item Only \bothAgnPct{}\% of galaxies are AGN by both BPT and WHAN.
        The \whanOnlyCount{} WHAN-only AGN are \woStrongPct{}\,\% Strong AGN; the \bptOnlyCount{}
        BPT-only AGN are \boRetPct{}\,\% Retired galaxies.
  \item \wiseMatched{} of \wiseAgnSample{} BPT AGN (\wiseMatchRate{}\%) have WISE counterparts;
        \wiseWedge{} of these (\wiseWedgePct{}\%) satisfy the Mateos+12 mid-infrared
        wedge, consistent with the wedge selecting only the most
        luminous AGN with strong hot dust emission.
  \item Electron densities from the [S\,II] ratio lie in
        $10$--$10^3$\,cm$^{-3}$ with no strong BPT trend.
  \item SMBH masses (\citet{tremaine2002}) peak at $\sim 3 \times 10^6\,M_\odot$,
        with \citet{gebhardt2000} and \citet{wu2009} giving consistent results.
\end{enumerate}

%---------------------------------------------------------------------
\appendix
\section{Tables}
\label{sec:appendix_tables}

Table~\ref{tab:sample_head} lists the first 10 galaxies in the sample
with their key measured and derived properties.

\begin{table}[htbp]
  \centering
  \caption{First 10 galaxies in the classified sample.}
  \label{tab:sample_head}
  \small
  \resizebox{\textwidth}{!}{%
  \begin{tabular}{rrrrrrcrr}
    \toprule
    specObjID & $z$ &
    $\log\,{\rm NII/H}\alpha$ & $\log\,{\rm OIII/H}\beta$ &
    EW(H$\alpha$) & BPT & WHAN &
    $\sigma_{[{\rm OIII}]}$ & $\log M_{\rm BH}$ \\
    \midrule
    299489677444933632 & 0.02 & $-0.32$ & $-0.24$ & 5.6 & Composite & wAGN & 124 & 7.29 \\
    29949050207\sigmaMedian{}54464 & 0.06 & $-0.28$ & $-0.08$ & 5.1 & Composite & wAGN & \whanPassCount{}6 & 7.17 \\
    299491051834468352 & 0.05 & $-0.69$ & $-0.03$ & 21.2 & SF & SF & 38 & 5.22 \\
    29949270\whanPassCount{}01910016 & 0.06 & $-0.50$ & $-0.42$ & 32.9 & SF & SF & 51 & 5.76 \\
    299493250857723904 & 0.06 & $-0.49$ & $-0.31$ & 24.8 & SF & SF & 76 & 6.43 \\
    299494075491444736 & 0.06 & $-0.44$ & $-0.51$ & 14.4 & SF & SF & 71 & 6.31 \\
    299496274514700288 & 0.13 & $-0.12$ & $-0.05$ & 4.8 & Composite & wAGN & 194 & 8.08 \\
    299496824270514176 & 0.01 & $-1.01$ & $0.12$ & 13.7 & SF & SF & 35 & 5.10 \\
    299499298171676672 & 0.05 & $-0.14$ & $-0.22$ & 15.4 & Composite & Strong AGN & 91 & 6.76 \\
    2995006725612\whanPassCount{}392 & 0.07 & $-0.38$ & $-0.57$ & 22.0 & SF & Strong AGN & \sigmaMedian{} & 6.65 \\
    \bottomrule
  \end{tabular}%
  }
\end{table}

\section{Python Code}
\label{sec:appendix_code}

The complete Python code is organised in the \texttt{agn\_project}
package and the \texttt{scripts/} directory.

\lstinputlisting[language=Python, caption=\texttt{agn\_project/config.py}]{../agn_project/config.py}
\lstinputlisting[language=Python, caption=\texttt{agn\_project/sdss.py}]{../agn_project/sdss.py}
\lstinputlisting[language=Python, caption=\texttt{agn\_project/classification.py}]{../agn_project/classification.py}
\lstinputlisting[language=Python, caption=\texttt{agn\_project/plots.py}]{../agn_project/plots.py}
\lstinputlisting[language=Python, caption=\texttt{agn\_project/wise.py}]{../agn_project/wise.py}
\lstinputlisting[language=Python, caption=\texttt{agn\_project/physics.py}]{../agn_project/physics.py}
\lstinputlisting[language=Python, caption=\texttt{agn\_project/utils.py}]{../agn_project/utils.py}

\bibliographystyle{apalike}
\bibliography{references}

\end{document}
